% Generated draft: requires mathematical and source review.
\documentclass[11pt]{article}
\usepackage[margin=25mm]{geometry}
\usepackage{fontspec}
\setmainfont{Latin Modern Roman}
\usepackage{amsmath,amssymb,mathtools,unicode-math}
\setmathfont{Latin Modern Math}
\usepackage{enumitem,needspace,xurl,hyperref}
\setlength{\parindent}{0pt}
\setlength{\parskip}{5pt}
\setlength{\emergencystretch}{3em}
\newcommand{\N}{\mathbb{N}}
\newcommand{\Z}{\mathbb{Z}}
\newcommand{\Q}{\mathbb{Q}}
\newcommand{\R}{\mathbb{R}}
\newcommand{\C}{\mathbb{C}}
\newcommand{\F}{\mathbb{F}}
\newcommand{\A}{\mathbb{A}}
\newcommand{\PP}{\mathbb{P}}
\newcommand{\E}{\mathbb{E}}
\newcommand{\eps}{\varepsilon}
\newcommand{\dd}{\,\mathrm{d}}
\begin{document}
\section*{UnsolvedMath ID: 30001269}
\textbf{Source problem code:} catalogue entry
\subsection{Definitions and mathematical statement}
Under the standard interpretation of the notation, let \(n\geq 0\) be an integer and let \(M\) be a closed smooth \(n\)-dimensional manifold equipped with a string structure \(S\). A string structure is a lift of the stable tangent classifying map from \(B\operatorname{Spin}\) to \(B\operatorname{String}\); equivalently, it includes a spin structure and a homotopy trivialization of its characteristic class \(\tfrac12p_1(TM)\). Let \(\operatorname{tmf}\) denote connective topological modular forms. Interpreting \(\sigma\) as the string orientation, write
\[
\sigma(M,S)\in\pi_n\operatorname{tmf}
\cong\operatorname{tmf}^{-n}(\mathrm{pt})
\]
for the image of the string bordism class of \((M,S)\).

The historical question in the supplied record is whether, for every such pair \((M,S)\),
\[
\left(\exists\text{ a Riemannian metric }g\text{ on }M\text{ with }
\operatorname{Ric}_g(v,v)>0
\text{ for every }x\in M\text{ and }0\ne v\in T_xM\right)
\quad\Longrightarrow\quad
\sigma(M,S)=0.
\]
The closed smooth domain, the meaning of \(\sigma\), and universal quantification over choices of \(S\) are standard interpretive conventions used to make this draft precise; the supplied statement does not explicitly specify them. It supplies no alternative domain, boundary conditions, or rule selecting \(S\).

\textbf{Source qualification.} The supplied original statement and verification note report that this implication has a negative answer in the cited workshop report. The verification note further says that a strengthened hypothesis is subsequently proposed, but does not supply that hypothesis. Thus this draft records the historical question and its supplied negative-answer qualification; it cannot state the strengthened question. The registry's open label and dated discussion of Stolz's related Witten-genus conjecture must be distinguished from the reported negative answer to this particular \(\operatorname{tmf}\)-valued question.
\subsection{Short English statement}
The historical question asks whether positive Ricci curvature forces the string orientation class of a string manifold to vanish in \(\operatorname{tmf}\). The supplied verification says this question has a negative answer and mentions an unspecified strengthened hypothesis. The record also labels the entry open while discussing Stolz's related Witten-genus conjecture. These qualifications concern different formulations and require human review.
\subsection{Sources}
\begin{itemize}
\item Supplied local UnsolvedMath registry record, permanent ID \(30001269\), title: Positive Ricci Curvature and String Orientation Classes; exact source problem code: \texttt{catalogue entry}. The record assigns the year 2009 and supplies no individual proposer.
\item \emph{Strings, Fields and Topology} (2009), Oberwolfach Reports, pp.\ 1547--1596, DOI: \texttt{10.4171/owr/2009/28}. Supplied source URL and literature-source link labeled original OWR report: \url{https://doi.org/10.4171/owr/2009/28}.
\item Supplied literature-source link labeled special cases of Stolz conjecture: \url{https://cris.uni-muenster.de/portal/en/publication/127470348}. The record does not supply its authors, publication title, or precise theorem.
\item Literature assessment dated exactly \texttt{2026-08-21}: the supplied curation says that the report records no counterexamples to Stolz's conjecture, that later work establishes special cases, and that no general resolution is known. This is supplied dated curation, not independently verified here. The background explicitly requests independent verification of time-sensitive status and cited claims.
\item The supplied statement status is \texttt{corrected\_verified}. Its verification note reports checking the official report and finding a negative answer to the naive positive-Ricci-curvature question, followed by a strengthened hypothesis. This is a supplied verification claim, not an independent source check performed for this draft. The registry nevertheless retains \texttt{open} status and an open-status tag.
\end{itemize}
\end{document}
